% Part: proof-theory
% Chapter: sequent-calculus
% Section: translations
%READERNOTE{SOL6-A337: G1c ते G3c भाषांतरातील उलट दिशेचा संदर्भ बदलला; दोन्ही दिशांत मुख्य संख्यापकीय सूत्र जपणाऱ्या नियमांचे आवश्यक दुर्बलीकरण किंवा आकुंचन स्पष्ट केले.}

\documentclass[../../../include/open-logic-section]{subfiles}

\begin{document}

\olfileid{pt}{seq}{tr}

\olsection{\Log{G1c} आणि \Log{G3c} यांतील रूपांतरे}

\Log{G1c} आणि \Log{G3c} या दोन्ही पद्धती अभिजात तर्कशास्त्रासाठी
निर्दोष व संपूर्ण आहेत, असे आपण म्हटले होते. म्हणजे प्रत्येक
!!{structure} मध्ये सत्य असलेल्या आणि फक्त त्याच क्रमवर्ती त्या
सिद्ध करतात. विशेषतः, दोन्ही त्याच क्रमवर्ती सिद्ध करतात. आता
निर्दोषता व संपूर्णतेच्या प्रमेयांचा वळसा न घालता, केवळ
सिद्धता-उपपत्तीय मार्गाने हे तथ्य सिद्ध करता येते. अशा
सिद्धतेत \Log{G1c} मधील !!{proof}s चे \Log{G3c} मधील
!!{proof}s मध्ये आणि उलट दिशेने \emph{रूपांतरे} मिळतात.

\begin{prop}
$\Log{G1c} \Proves S$ असेल, तर $\Log{G3c} \Proves S$.
\end{prop}

\begin{proof}
  $\pi$ ही \Log{G1c} मधील !!a{proof} असेल, तर \Log{G3c} मध्ये
  $S$ ची !!a{proof}~$\pi'$ आहे, हे दाखवू. यासाठी $n
  = \pheight{\pi}$ वर विगमन करू.

  $n = 0$ असेल, तर $\pi$ हे स्वयंसिद्धक असते. $\lfalse \Sequent
  \quad$ हे स्वयंसिद्धक \Log{G3c} मध्येही आहे (संदर्भ $\Gamma,
  \Delta$ रिकामे घ्या). \Log{G1c} मधील $!A \Sequent !A$ ही
  स्वयंसिद्धके $!A$ अण्वीय असण्याची अट न घालता ग्राह्य असतात.
  पण अशा सर्व क्रमवर्ती \Log{G3c} मध्ये !!{provable} आहेत, हे आपण
  \olref[ex]{prop:G3c-axioms} मध्ये सिद्ध केले आहे.

  $n>0$ असेल, तर $\pi$ चा शेवट $R$ नियमाने होतो. विगमनाचे
  गृहीतक सांगते की त्या नियमाच्या आधारविधानांच्या \Log{G3c}
  मधील !!{proof}s आहेत; म्हणजे लगेचच्या उप-!!{proof}s ची
  \Log{G3c} मध्ये रूपांतरे आहेत. $R$ हा \Log{G3c} मधीलही नियम
  असेल, तर आधारविधानांच्या रूपांतरित !!{proof}s वर तो लागू करू.
  $R$ हा \Log{G1c} मधील $\LeftR{\land}$ असेल,
  तर आधारविधानात नसलेले दुसरे संयोजक सूत्र डावीकडे
  दुर्बलीकरणाने जोडा आणि \Log{G3c} चा $\LeftR{\land}$
  नियम लावा. $\RightR{\lor}$ साठी नसलेले दुसरे वियोजक
  सूत्र उजवीकडे जोडा आणि \Log{G3c} चा $\RightR{\lor}$
  नियम लावा. \Log{G1c} च्या $\LeftR{\lforall}$ किंवा
  $\RightR{\lexists}$ साठी त्याच बाजूला मुख्य संख्यापकीय
  सूत्र दुर्बलीकरणाने आधी जोडा; मग ते जपणारा \Log{G3c}
  नियम लावा. दुर्बलीकरणाची अनुज्ञेयता
  \olref[adm]{prop:weak-G3c-adm} देते. $R$ स्वतः
  दुर्बलीकरणाचा नियम असेल, तर तीच प्रमेयिका वापरा;
  आकुंचनाचा नियम असेल, तर
  \olref[inv]{prop:G3c-cont-adm} वापरा.
\end{proof}

\begin{prop}
$\Log{G3c} \Proves S$ असेल, तर $\Log{G1c} \Proves S$.
\end{prop}

\begin{proof}
  पुन्हा $S$ ची !!a{proof}~$\pi$ हिच्या !!{height} वर विगमन करू.
  \Log{G3c} मधील प्रत्येक स्वयंसिद्धक \Log{G1c} मध्ये एका
  स्वयंसिद्धकापासून \LeftR{\Weakening} आणि
  \RightR{\Weakening} नियम काही वेळा वापरून !!{prove}d करता येते:
  \[
  \Axiom$!C \fCenter !C$
  \RightLabel{\LeftR{\Weakening}}
  \Deduce$!C, \Gamma \fCenter !C$
  \RightLabel{\RightR{\Weakening}}
  \Deduce$!C, \Gamma \fCenter \Delta, !C$
  \DisplayProof
\qquad
  \Axiom$\lfalse \fCenter$
  \RightLabel{\LeftR{\Weakening}}
  \Deduce$\lfalse, \Gamma \fCenter$
  \RightLabel{\RightR{\Weakening}}
  \Deduce$\lfalse, \Gamma \fCenter \Delta$
  \DisplayProof
  \]
  $\pi$ चा शेवट $R$ नियमाने होत असेल, तर विगमनाच्या गृहीतकाने
  आधारविधानांच्या \Log{G1c} मधील !!{proof}s मिळतात;
  नियम दोन्ही पद्धतींत समान असेल, तर त्या !!{proof}s वर
  तोच $R$ नियम लागू करता येतो. $\LeftR{\lforall}$ आणि
  $\RightR{\lexists}$ मध्ये \Log{G3c} च्या आधारविधानात
  मुख्य संख्यापकीय सूत्रही जपलेले असते. ते संदर्भात ठेवून
  \Log{G1c} चा नियम लावल्यास निष्कर्षात मुख्य सूत्राच्या
  दोन प्रती मिळतात; अनुक्रमे डाव्या किंवा उजव्या आकुंचनाने
  त्या एक करता येतात. उरलेले अपवाद म्हणजे \LeftR{\land}
  आणि \RightR{\lor}; त्यांच्या \Log{G1c} आणि
  \Log{G3c} मधील आवृत्त्या वेगळ्या आहेत. \Log{G3c} मधील
  आवृत्त्या \Log{G1c} मध्ये पुढीलप्रमाणे व्युत्पाद्य आहेत:
  \[
  \Axiom$!A, !B, \Gamma \fCenter \Delta$
  \RightLabel{\LeftR{\land}}
  \UnaryInf$!A \land !B, !B, \Gamma \fCenter \Delta$
  \RightLabel{\LeftR{\land}}
  \UnaryInf$!A \land !B, !A \land !B, \Gamma \fCenter \Delta$
  \RightLabel{\LeftR{\Contraction}}
  \UnaryInf$!A \land !B, \Gamma \fCenter \Delta$
    \DisplayProof
\qquad
  \Axiom$\Gamma \fCenter \Delta, !A, !B$
  \RightLabel{\RightR{\lor}}
  \UnaryInf$\Gamma \fCenter \Delta, !A \lor !B, !B$
  \RightLabel{\RightR{\lor}}
  \UnaryInf$\Gamma \fCenter \Delta, !A \lor !B, !A \lor !B$
  \RightLabel{\RightR{\Contraction}}
  \UnaryInf$\Gamma \fCenter \Delta, !A \lor !B$
    \DisplayProof
  \]
\end{proof}

\end{document}
